% ============================================================================
\section{Discussion et limites}
\label{sec:limites}
% ============================================================================

\subsection{Ce que les mesures nous ont appris}
\begin{enumerate}
  \item \textbf{Mesurer avant d'optimiser change le problème.} Deux idées
        intuitives se sont révélées fausses : HighCPU n'est pas plus rapide
        par tâche, et une VM à 2~vCPU n'exécute pas 2 tâches à pleine vitesse.
        En tenir compte a fait passer l'erreur du modèle de 65\,\% à 8\,\%
        (tableau~\ref{tab:validation}), et a montré que le vrai levier est le
        choix des VM à allumer.
  \item \textbf{Le choix des VM compte plus que la répartition des tâches},
        surtout pour les petits lots, à cause de la facturation minimale de
        60~s.
  \item \textbf{Une règle fixe ne suffit pas} : éteindre vm1 est le bon choix
        sur 500 tâches, mais pas sur 2000. Seul le calcul de toutes les
        combinaisons permet de trancher.
\end{enumerate}

\subsection{Limites}

\paragraph{Variabilité des VM à cœur partagé.} C'est la limite principale. Les
e2-medium sont les VM les moins chères, et donc celles que préfère notre
algorithme. Mais GCP ne leur garantit qu'un vCPU sur deux : au-delà, elles
dépendent de rafales limitées, et leur vitesse peut chuter de moitié sous charge
prolongée (section~\ref{sec:vm3}). Notre modèle, calibré sur des mesures courtes
(tableau~\ref{tab:debitprix}), ne le prévoit pas, et un planning calculé à
l'avance ne s'y adapte pas.

Nous avons simulé ce pire cas en modélisant les deux e2-medium à mi-vitesse
($\gamma = 2{,}8$ au lieu de 1,4). \sel{} change alors de choix : highcpu seule
sur 9 tâches, une seule e2-medium et highcpu sur 500 tâches, toujours les 4 VM
sur 2000 tâches. Sur 9 tâches, ce scénario ne se produit pas en pratique : aucune
VM n'a été bridée sur un lot aussi court. Sur les lots longs, un modèle à deux
vitesses (pleine vitesse pendant la rafale, puis vitesse garantie) donnerait des
choix plus prudents.

\paragraph{Planning statique.} Le scheduler décide tout avant l'exécution ; si une
VM ralentit, ses tâches attendent même si d'autres VM ont fini.

\paragraph{Peu de répétitions.} Chaque algorithme n'a été exécuté que 3 fois par
session. C'est suffisant pour les écarts de coût, stables d'une exécution à
l'autre, mais pas pour départager des makespans proches quand une VM est bridée.

\paragraph{Contention appliquée partout.} Le modèle ralentit toutes les tâches
de $\times 1{,}4$, comme si chaque VM était pleinement occupée. C'est juste sur les
gros lots, mais trop pessimiste sur les petits, où une VM peut n'exécuter qu'une
tâche : sur 9 tâches, Min-Min (une seule tâche sur vm2), prédit le plus lent,
s'est révélé le plus rapide.

\paragraph{Contrainte de délai relative.} L'échéance $D \le 1{,}25 \cdot \Drr$
dépend de RR : plus RR est lent, plus elle est lâche. Elle sert à comparer les
algorithmes sur des lots de tailles très différentes ; pour un client réel, on la
remplacerait par son échéance contractuelle, sans changer l'algorithme.

\paragraph{Coût calculé, pas facture réelle.} Notre coût ne compte que le temps
de travail des VM (au moins 60~s), sans le démarrage, les temps morts, le
disque ni l'IP. Il sert à \emph{comparer} les allocations, pas à chiffrer une
économie réelle.

\subsection{Pistes d'amélioration}

\begin{itemize}
  \item \textbf{Modéliser deux vitesses pour les e2-medium} : pleine vitesse
        pendant la rafale, environ moitié au-delà. Les prédictions des lots longs
        seraient plus prudentes et \sel{} enverrait moins de tâches aux e2-medium.
  \item \textbf{Mesurer la vitesse réelle des VM juste avant de planifier} (quelques
        tâches de test), et mettre à jour le modèle.
  \item \textbf{Rééquilibrage dynamique} (\emph{work stealing}) : une VM qui a fini
        prend des tâches restantes d'une VM lente, au prix d'un écart au planning.
  \item \textbf{Surveiller le \emph{steal time}} des VM pour confirmer et quantifier
        le bridage des e2-medium.
  \item \textbf{Recherche locale} après \sel{} (déplacements et échanges qui font
        baisser $S$), pour affiner la répartition au-delà de Max-Min.
  \item \textbf{Plus de répétitions} (10 par algorithme) pour départager les
        makespans proches.
\end{itemize}
